\begin{abstract}
Contemporary language models can generate fluent responses that are factually incorrect or unfaithful to the input context. We propose HIDE, a training-free, single-generation hallucination detector motivated by input--output representational dependence. HIDE applies an adapted HSIC-based score to selected hidden states collected during generation. We evaluate six open-source language models on four QA benchmarks covering faithfulness and factuality. In the original benchmark tables, HIDE has approximately $29\%$ mean relative AUC improvement over the better of the evaluated perplexity and energy baselines within each model/dataset setting, and is competitive with multi-generation methods. A new matched comparison on 2,000 examples per dataset and model, using explicit first-answer-line stopping, finds that prompt attention mass outperforms HIDE in seven of eight settings. A selected-token-count control closely tracks HIDE, limiting a specifically nonlinear-dependence interpretation of its performance. Paired latency measurements extend to Gemma-2-27B and show model- and dataset-dependent overhead under standard HuggingFace inference. The historical approximately $51\%$ reduction relative to the evaluated multi-pass methods is a baseline architectural comparison, not an optimized-serving measurement. These results identify both the practical scope and limitations of the proposed score.
\footnote{Code and retained experimental records: \url{https://github.com/C-anwoy/HIDE}.}
\end{abstract}
